% Provenance: copied from the design workspace CambridgePQ/recursion.fast, `report/sections/01-summary.tex`, on 2026-10-19. Read-only copy; the source is authoritative.
\section{Executive summary}\label{sec:summary}

\paragraph{The problem.} recursion.fast asks agents to design SNARK verifiers that are as cheap as
possible to recurse on. The design space should stay wide open, while a small trusted base decides
what counts as a correct verifier and what it costs. The earlier Lean-checked competitions show the
format works: ots.golf, sig.golf and the Proximity Prize (\cref{sec:prior}). They also show its
failure modes: vacuous statements, images unbound from theorems, flat per-instruction cost models
that agents overfit, and check times near their limits. This report turns the brief into a
specification with those lessons applied. Every decision below was attacked by at least one
independent adversarial review, and the load-bearing parts were built and checked in Lean.

\begin{takeaway}[Headline recommendations]
\begin{enumerate}[nosep,leftmargin=*]
\item \textbf{Objective.} Score the cost of one recursion step: a proved bound on the cost of every
  \emph{accepting} run of the submitted verifier image (\cref{sec:objectives,sec:cost}). Rejecting
  runs are fuel-capped and unscored, so fast-failing paths are free. Proof size and prover cost are
  constraints and secondary boards. The brief's \enquote{performance $\div$ program size} is dropped
  (\cref{sec:brief-review}).
\item \textbf{Model.} A holographic public-coin IOP in the information-theoretic model
  (\cref{sec:protocol}). R1CS over $K=\mathrm{GF}(2^{64})$ with challenges in
  $E=\mathrm{GF}(2^{192})$, at three fixed size classes $2^{16},2^{20},2^{24}$ with
  satisfiability-preserving and reflecting padding. Prover messages are raw $K$-vectors; the
  challenge menu and a round cap of $2^{14}$ are fixed by the contract. BCS and Fiat--Shamir
  compilation is a cited meta-theorem: per-round error $2^{-128}$ gives a soundness error of at most
  $t/2^{127}$ for $2k\le t\le2^{124}$, which is a knowledge error only with unbounded extraction time
  (\cref{subsec:sec-cap}).
\item \textbf{Security definitions.} Perfect completeness and one worst-case round-by-round
  knowledge-soundness notion in counting form, restated locally in 132 lines and bridged to
  ArkLib. Because Lean extractors carry no cost bound, level-1 records are labelled
  \enquote{RBR-sound}. Knowledge claims require a separate efficient-extraction track
  (\cref{subsec:sec-extractor}).
\item \textbf{Binding.} The RISC-V image must agree exactly with a Lean spec verifier on every raw
  transcript, every index of class shape and every instance, with honest preprocessing
  (\cref{subsec:proto-agrees}). Images and score inputs are rendered from files into the challenge
  statement, and no candidate code runs on the trust path (\cref{sec:architecture}).
\item \textbf{Cost.} One proved 18-counter vector per class. The \emph{recursion board} prices it
  under four calibrated zkVM profiles (ZisK, SP1-Hypercube, RISC Zero, Jolt) and takes the
  geometric mean. A \emph{latency board} uses a deterministic in-order timing model and fork--join
  harts, with work and span accounting. Wall time on a reference board with Zbc is used only for
  calibration. This answers the overfitting seen under sig.golf's flat per-instruction model
  (\cref{sec:cost}).
\item \textbf{Anti-overfitting.} The cost bound is a class-uniform theorem, over \emph{every} R1CS
  of the class, so hard-coding benchmark circuits buys nothing. Benchmarks are used for
  demonstration and the record gate (\cref{sec:benchmark}).
\item \textbf{Recursion.} The batcher is an accumulation scheme (embed, fold, decide) with a fixed,
  typed accumulator shape (\cref{sec:recursion}). One-step theorems are required at level~1 and
  the tree theorem is an organiser obligation. Soundness of the deployed recursive SNARK is the named
  assumption \lean{RecursionHeuristic}, which no Lean theorem can currently replace.
\item \textbf{MVP.} The trivial verifier (proof = $w$) is fully certified against the final
  contract v0.3, including the optional latency clause \lean{Latency1}. The proof is 1{,}920 lines.
  The production harness accepts it with two kernels, at about 6 minutes per harness run and
  3.5\,GB; the full check with differential tests takes 13.6 minutes. It scores
  $\mathrm{Score}_{\mathrm{zk}}\approx1.6\times10^9$ and $\mathrm{Score}_{\mathrm{lat}}\approx6.1\times10^8$,
  and its cost bounds are tight (\cref{sec:baseline}). Modelled succinct designs are about two
  orders of magnitude better (model estimates, accurate to about $2\times$).
\end{enumerate}
\end{takeaway}

\paragraph{Answers to the brief's five questions.}
\begin{description}[leftmargin=0pt,itemsep=2pt]
\item[Objective and design.] The aim is general, cheap-to-recurse verifier architectures, with
  security certified in Lean and cost certified on the exact bytes that would be deployed. Agents
  control the protocol, the indexer, the honest prover, the RISC-V image and the proofs. The
  organisers control a trusted contract (796 Lean code lines in v0.3, over 930 reachable lines of
  vendored RISC-V semantics), a harness, and an untrusted organiser library
  (\cref{sec:objectives,sec:architecture}).
\item[Rules, trust and gaps.] The harness renders the challenge, then runs comparator first, in a
  fresh tree, with an axiom whitelist and two kernel replays, then an environment-level hygiene gate
  in a second fresh tree, with the organiser build pinned by a hash manifest. Records also pass an
  unscored badge path, which v0.3 specifies but does not yet implement (\cref{sec:architecture}).
  What must be trusted is listed explicitly: the contract, the Lean kernel, elaborator, Lake and
  compiler, comparator and the second kernel, the vendored RISC-V semantics, the harness and its host
  tools, the cited meta-theorems, any named hypotheses, and the deployment assumptions. The remaining
  gaps are in \cref{sec:threats,sec:issues}.
\item[Submission format and cost.] Raw RV64IM+Zbc images, a claim file, Lean proofs, and native
  binaries for record runs. The cost is the proved accepting-run counter vector priced by calibrated
  profiles, plus a Fiat--Shamir surcharge, a byte charge and a soft prover blow-up factor. Every
  attempted gaming vector is catalogued with its mitigation (\cref{sec:cost,sec:threats}).
\item[Security definitions.] Perfect completeness, worst-case RBR knowledge soundness with a named
  extractor, the per-round cap $2^{-128}$, the round cap $2^{14}$, and controls showing every clause
  is load-bearing (\cref{sec:security}).
\item[Baseline.] Certified and timed as above, first against v0.2 (1{,}481 lines, 2.5 minutes on
  cached dependencies) and then against v0.3. The brief's one-hour requirement is met with a margin of
  more than $4\times$ for the complete v0.3 check (\cref{sec:baseline}).
\end{description}

\paragraph{Answers to the brief's further considerations.}
\begin{description}[leftmargin=0pt,itemsep=2pt]
\item[MVP scope.] Launch in phases (\cref{subsec:iss-roadmap}): freeze contract v0.3 with the
  certified trivial baseline (done) and organiser library components lib-L0 and lib-L1; open a
  Lean-only T0 season, in which a verifier is a query program scored on its three query counters;
  then open T1, RISC-V images on the recursion board. The latency board, fork--join harts and the
  accumulation track (T2) are optional or later.
\item[An easy-to-implement cost function.] The score is a Lean definition of the contract
  (\code{Cost.lean}, 73 code lines including the cost clauses, the four profiles, the score
  functions and the admission caps); the external scorer only takes a root of a certified rational.
  A submitter proves one counter vector per class, and re-weighting the profiles needs no new proof
  (\cref{subsec:arch-versioning}). It is not the simplest possible model: a single flat weight
  vector, kept as a display column, would be simpler, and was rejected because agents overfit it
  (\cref{subsec:cost-diagnosis}). A minimal launch can omit the latency board and harts.
\item[Design-space exploration.] The contract fixes the relation, fields, machine and round menu,
  and nothing about the protocol (\cref{subsec:obj-interesting}); an untrusted organiser library
  lowers the proof burden so that it does not select designs.
\item[Inner versus final verifier.] The batcher is an accumulation scheme; the fold verifier (the
  inner verifier) is scored at every node and the decider (the final verifier) at weight $1/64$
  (\cref{subsec:rec-score}).
\item[Interactive or non-interactive.] Interactive: submitters prove properties of a public-coin IOP
  in the information-theoretic model, and compilation to a non-interactive argument is a cited
  theorem with an explicit bound, plus a harness rule for the Fiat--Shamir transcript (Decision~P-1,
  \cref{subsec:sec-fs}).
\item[Overfitting.] Knowledge soundness cannot prevent it, so the cost theorem is class-uniform
  (\cref{subsec:bench-hardcode}); overfitting of the cost model itself is answered by several
  calibrated profiles (\cref{subsec:cost-diagnosis}).
\end{description}

\paragraph{User inputs during the study.} Three inputs from the brief's author were incorporated. The
competition is over binary fields (U1): \cref{subsec:bg-binary,subsec:proto-fields}. The cost bound is
taken over accepting runs, with a fast-failing path (U2): \cref{sec:cost}. Flat per-instruction costs
are replaced by calibrated and timing models with parallelism (U3): \cref{sec:cost}.

\paragraph{Evidence produced.} The study produced:
\begin{itemize}[nosep]
\item eight research notes and four design documents;
\item eight adversarial reviews: incentives, definitions, feasibility, the compiled contract in two
  rounds (v0.2 and v0.3), the cryptography, the formal methods and the editing of the report;
\item a compiled trusted contract: v0.2 with 582 code lines, and v0.3 with 796, which adds the
  cost-realism and review fixes;
\item Lean counterexamples for every critical definitional attack found, each now a failing control;
\item a certified baseline submission;
\item cycle-level probes of field kernels under nine cost models.
\end{itemize}
All artefacts are reproducible from the paths in \cref{sec:reproducibility}.

\paragraph{Main risks.} The decisive risk is proof burden. A competitive entry is sig.golf-scale work,
estimated at 45--100k lines with organiser libraries. Its critical path is a knowledge-sound
hash-based PCS over binary Reed--Solomon codes, which neither ArkLib nor any other Lean library
provides (\cref{sec:issues}). Further risks:
\begin{itemize}[nosep]
\item the Johnson-regime proximity gap, available only as a flagged named hypothesis;
\item recursive soundness after hash instantiation;
\item the calibration of judgement constants (profiles, surcharge, blow-up threshold);
\item deployment fidelity of the per-hart fuel rule, and three score-meaning gaps (hart pricing,
  timer realism, leaf pricing on the latency board) scheduled for v0.4;
\item harness components that are specified but not yet built (production sandbox, badge path,
  trusted emulator, Fiat--Shamir sampler);
\item contract churn after launch.
\end{itemize}
The roadmap sequences these: contract freeze and library components lib-L0--lib-L1, then a Lean-only
T0 season, then T1 with the PCS library, then T2 recursion (\cref{subsec:iss-roadmap}).

\paragraph{How to read this report.} Cryptographers new to Lean should read
\cref{subsec:bg-lean} before \cref{sec:protocol,sec:security}. Formal-methods readers new to SNARKs
should read \cref{subsec:bg-iop,subsec:bg-bcs,subsec:bg-recursion}. zkVM engineers will find the
machine and cost model in \cref{subsec:proto-machine,sec:cost}. Decisions appear in blue boxes,
actionable issues in red boxes and trusted assumptions in orange boxes. A glossary is in
\cref{sec:glossary}. Four families of short names recur: T0--T3 are score tracks; L0--L4 are
security levels (\cref{subsec:proto-model}); lib-L0--lib-L6 are organiser library components
(\cref{sec:issues}); and the latency-board record is named after its contract clause
\lean{Latency1} (its challenge stub is the file \code{L1.lean.in}, which is not security level L1).
Bracketed tags name the source of a rule: A1--A8 are the adversarial reviews (A1 incentives, A2
definitions, A3 feasibility, A4 contract v0.2, A5 cryptography of the report, A6 contract v0.3, A7
formal methods of the report, A8 editing); D1--D4 are design notes and \enquote{DESIGN §$n$} the
design document; P2--P3b are engineering reports (contract, baseline); O-recheck is an independent
re-run of the baseline's checks. These documents ship with the report.
